\documentclass[12pt,a4paper]{article}

% ------------------------------------------------------------
% Кодировка и русский язык
% ------------------------------------------------------------
\usepackage[T2A]{fontenc}
\usepackage[utf8]{inputenc}
\usepackage[english,russian]{babel}
\usepackage{cmap}

% ------------------------------------------------------------
% Математика
% ------------------------------------------------------------
\usepackage{amsmath,amssymb,amsthm,mathtools}
\usepackage{bm}

% ------------------------------------------------------------
% Оформление
% ------------------------------------------------------------
\usepackage{geometry}
\geometry{
    left=25mm,
    right=20mm,
    top=20mm,
    bottom=20mm
}

\usepackage[expansion=false]{microtype}
\usepackage{xcolor}
\usepackage{enumitem}
\usepackage{booktabs}
\usepackage{array}
\usepackage{tabularx}
\usepackage{longtable}
\usepackage{multirow}
\usepackage{tikz}
\usepackage[edges]{forest}
\usepackage[most]{tcolorbox}
\usepackage{hyperref}

\hypersetup{
    unicode=true,
    colorlinks=false,
    pdfborder={0 0 0}
}

\setlength{\parindent}{1.25cm}
\setlength{\parskip}{0.25em}
\setlength{\emergencystretch}{3em}

% ------------------------------------------------------------
% Теоремы и определения
% ------------------------------------------------------------
\theoremstyle{plain}
\newtheorem{theorem}{Теорема}[section]
\newtheorem{proposition}[theorem]{Утверждение}
\newtheorem{lemma}[theorem]{Лемма}
\newtheorem{corollary}[theorem]{Следствие}

\theoremstyle{definition}
\newtheorem{definition}[theorem]{Определение}
\newtheorem{example}[theorem]{Пример}

\theoremstyle{remark}
\newtheorem{remark}[theorem]{Замечание}

\renewcommand{\proofname}{Доказательство}

% ------------------------------------------------------------
% Пользовательские команды
% ------------------------------------------------------------
\newcommand{\Pp}{\mathbb{P}}
\newcommand{\Ee}{\mathbb{E}}
\newcommand{\Rr}{\mathbb{R}}
\newcommand{\ind}{\mathbf{1}}
\newcommand{\argmax}{\operatorname*{arg\,max}}
\newcommand{\argmin}{\operatorname*{arg\,min}}
\newcommand{\Entropy}{\operatorname{H}}
\newcommand{\Gain}{\operatorname{Gain}}
\newcommand{\Gini}{\operatorname{Gini}}
\newcommand{\Risk}{\operatorname{R}}

% ------------------------------------------------------------
% Блоки алгоритмов и важных замечаний
% ------------------------------------------------------------
\newtcolorbox{algorithmblock}[1]{
    breakable,
    enhanced,
    colback=blue!3,
    colframe=blue!45!black,
    fonttitle=\bfseries,
    title={#1},
    boxrule=0.7pt,
    arc=2mm
}

\newtcolorbox{importantblock}[1]{
    breakable,
    enhanced,
    colback=yellow!5,
    colframe=orange!55!black,
    fonttitle=\bfseries,
    title={#1},
    boxrule=0.7pt,
    arc=2mm
}

\newtcolorbox{exampleblock}[1]{
    breakable,
    enhanced,
    colback=green!3,
    colframe=green!45!black,
    fonttitle=\bfseries,
    title={#1},
    boxrule=0.7pt,
    arc=2mm
}

% ------------------------------------------------------------

\title{Билет №95 \\ \small Теорема Байеса. Деревья решений. (ИТМО)}
\author{Сериков Владислав Александрович}
\date{19 июля 2026}

\begin{document}

\maketitle
\tableofcontents
\newpage

% ============================================================
\section{Условная вероятность и теорема Байеса}
% ============================================================

\subsection{Вероятностное пространство и условная вероятность}

Пусть
\[
    (\Omega,\mathcal{F},\Pp)
\]
--- вероятностное пространство, где:
\begin{itemize}
    \item $\Omega$ --- множество элементарных исходов;
    \item $\mathcal{F}$ --- $\sigma$-алгебра событий;
    \item $\Pp$ --- вероятностная мера.
\end{itemize}

\begin{definition}
Пусть $A,B\in\mathcal{F}$ и $\Pp(B)>0$. \emph{Условной вероятностью}
события $A$ при условии наступления события $B$ называется число
\[
    \Pp(A\mid B)
    =
    \frac{\Pp(A\cap B)}{\Pp(B)}.
\]
\end{definition}

Условная вероятность показывает, как изменяется вероятность события $A$,
если известно, что событие $B$ уже произошло.

Из определения непосредственно получается \emph{формула умножения}:
\[
    \Pp(A\cap B)
    =
    \Pp(A\mid B)\Pp(B).
\]
Если также $\Pp(A)>0$, то
\[
    \Pp(A\cap B)
    =
    \Pp(B\mid A)\Pp(A).
\]
Следовательно,
\[
    \Pp(A\mid B)\Pp(B)
    =
    \Pp(B\mid A)\Pp(A).
\]

\begin{definition}
События $A$ и $B$ называются \emph{независимыми}, если
\[
    \Pp(A\cap B)=\Pp(A)\Pp(B).
\]
Если $\Pp(B)>0$, это условие эквивалентно равенству
\[
    \Pp(A\mid B)=\Pp(A).
\]
\end{definition}

Таким образом, при независимости информация о наступлении $B$ не меняет
вероятность события $A$.

\subsection{Разбиение пространства событий}

\begin{definition}
События $H_1,\ldots,H_m$ образуют \emph{полную группу событий}, или
\emph{разбиение} пространства $\Omega$, если:
\[
    H_i\cap H_j=\varnothing,\qquad i\ne j,
\]
и
\[
    \bigcup_{i=1}^{m}H_i=\Omega.
\]
\end{definition}

События $H_i$ часто интерпретируются как взаимно исключающие гипотезы.
Если заранее известна вероятность $\Pp(H_i)$, то она называется
\emph{априорной вероятностью} гипотезы $H_i$.

\subsection{Формула полной вероятности}

\begin{theorem}[Формула полной вероятности]
Пусть события $H_1,\ldots,H_m$ образуют разбиение пространства $\Omega$,
причём $\Pp(H_i)>0$ для всех $i$. Тогда для любого события $A$ справедливо
равенство
\[
    \Pp(A)
    =
    \sum_{i=1}^{m}
    \Pp(A\mid H_i)\Pp(H_i).
\]
\end{theorem}

\begin{proof}
Поскольку события $H_1,\ldots,H_m$ образуют разбиение пространства,
имеем
\[
    \Omega=\bigcup_{i=1}^{m}H_i.
\]
Следовательно,
\[
    A=A\cap\Omega
     =A\cap\left(\bigcup_{i=1}^{m}H_i\right)
     =\bigcup_{i=1}^{m}(A\cap H_i).
\]
События $A\cap H_i$ попарно несовместны, поскольку попарно несовместны
события $H_i$. По конечной аддитивности вероятности
\[
    \Pp(A)
    =
    \sum_{i=1}^{m}\Pp(A\cap H_i).
\]
По формуле умножения
\[
    \Pp(A\cap H_i)
    =
    \Pp(A\mid H_i)\Pp(H_i).
\]
Подставляя эти равенства в предыдущую сумму, получаем
\[
    \Pp(A)
    =
    \sum_{i=1}^{m}
    \Pp(A\mid H_i)\Pp(H_i).
\]
\end{proof}

\subsection{Теорема Байеса}

\begin{theorem}[Теорема Байеса]
Пусть $H_1,\ldots,H_m$ образуют разбиение пространства $\Omega$, причём
$\Pp(H_i)>0$ для всех $i$. Пусть $A$ --- событие, для которого
$\Pp(A)>0$. Тогда для каждого $k\in\{1,\ldots,m\}$
\[
    \boxed{
    \Pp(H_k\mid A)
    =
    \frac{\Pp(A\mid H_k)\Pp(H_k)}
    {\displaystyle
        \sum_{i=1}^{m}\Pp(A\mid H_i)\Pp(H_i)}
    }.
\]
\end{theorem}

\begin{proof}
По определению условной вероятности
\[
    \Pp(H_k\mid A)
    =
    \frac{\Pp(H_k\cap A)}{\Pp(A)}.
\]
События $H_k\cap A$ и $A\cap H_k$ совпадают. По формуле умножения
\[
    \Pp(H_k\cap A)
    =
    \Pp(A\mid H_k)\Pp(H_k).
\]
Следовательно,
\[
    \Pp(H_k\mid A)
    =
    \frac{\Pp(A\mid H_k)\Pp(H_k)}
         {\Pp(A)}.
\]
По формуле полной вероятности
\[
    \Pp(A)
    =
    \sum_{i=1}^{m}
    \Pp(A\mid H_i)\Pp(H_i).
\]
Подставляя это выражение в знаменатель, получаем
\[
    \Pp(H_k\mid A)
    =
    \frac{\Pp(A\mid H_k)\Pp(H_k)}
    {\displaystyle
        \sum_{i=1}^{m}\Pp(A\mid H_i)\Pp(H_i)}.
\]
\end{proof}

Для двух событий $A$ и $B$ теорема Байеса имеет вид
\[
    \boxed{
    \Pp(A\mid B)
    =
    \frac{\Pp(B\mid A)\Pp(A)}{\Pp(B)}
    },
    \qquad \Pp(B)>0.
\]

Если в качестве разбиения используются события $A$ и $\overline{A}$, то
\[
    \Pp(A\mid B)
    =
    \frac{\Pp(B\mid A)\Pp(A)}
    {\Pp(B\mid A)\Pp(A)
    +
    \Pp(B\mid \overline{A})\Pp(\overline{A})}.
\]

\subsection{Интерпретация формулы Байеса}

Формулу Байеса часто записывают в виде
\[
    \underbrace{\Pp(H\mid A)}_{\text{апостериорная вероятность}}
    =
    \frac{
        \underbrace{\Pp(A\mid H)}_{\text{правдоподобие}}
        \underbrace{\Pp(H)}_{\text{априорная вероятность}}
    }{
        \underbrace{\Pp(A)}_{\text{свидетельство}}
    }.
\]

Здесь:
\begin{itemize}
    \item $\Pp(H)$ --- степень уверенности в гипотезе до получения новых
    данных;
    \item $\Pp(A\mid H)$ --- вероятность наблюдать данные $A$, если
    гипотеза $H$ верна;
    \item $\Pp(A)$ --- полная вероятность наблюдаемых данных;
    \item $\Pp(H\mid A)$ --- обновлённая вероятность гипотезы после
    наблюдения данных.
\end{itemize}

В знаменателе находится нормирующая константа, обеспечивающая равенство
\[
    \sum_{i=1}^{m}\Pp(H_i\mid A)=1.
\]

\subsection{Форма в терминах отношений шансов}

Для гипотез $H_1$ и $H_0$ теорема Байеса даёт
\[
    \frac{\Pp(H_1\mid A)}{\Pp(H_0\mid A)}
    =
    \frac{\Pp(A\mid H_1)}{\Pp(A\mid H_0)}
    \cdot
    \frac{\Pp(H_1)}{\Pp(H_0)}.
\]

Иными словами,
\[
    \boxed{
    \text{апостериорные шансы}
    =
    \text{отношение правдоподобия}
    \times
    \text{априорные шансы}
    }.
\]

Величина
\[
    \frac{\Pp(A\mid H_1)}{\Pp(A\mid H_0)}
\]
называется \emph{отношением правдоподобия}.

Если последовательно наблюдаются условно независимые данные
$A_1,\ldots,A_n$, то
\[
    \frac{\Pp(H_1\mid A_1,\ldots,A_n)}
         {\Pp(H_0\mid A_1,\ldots,A_n)}
    =
    \frac{\Pp(H_1)}{\Pp(H_0)}
    \prod_{j=1}^{n}
    \frac{\Pp(A_j\mid H_1)}{\Pp(A_j\mid H_0)}.
\]

\subsection{Непрерывная форма теоремы Байеса}

Пусть $\Theta$ --- неизвестный параметр с априорной плотностью
$\pi(\theta)$, а $X$ --- наблюдаемая случайная величина с условной
плотностью $p(x\mid\theta)$. Совместная плотность имеет вид
\[
    p(x,\theta)
    =
    p(x\mid\theta)\pi(\theta).
\]

Маргинальная плотность наблюдения равна
\[
    p(x)
    =
    \int p(x\mid\vartheta)\pi(\vartheta)\,d\vartheta.
\]

Если $p(x)>0$, то апостериорная плотность параметра равна
\[
    \boxed{
    \pi(\theta\mid x)
    =
    \frac{
        p(x\mid\theta)\pi(\theta)
    }{
        \displaystyle
        \int p(x\mid\vartheta)\pi(\vartheta)\,d\vartheta
    }
    }.
\]

Часто используют запись с точностью до нормирующей константы:
\[
    \pi(\theta\mid x)
    \propto
    p(x\mid\theta)\pi(\theta).
\]

При фиксированных данных $x$ функция
\[
    L(\theta;x)=p(x\mid\theta)
\]
называется \emph{функцией правдоподобия}. Важно, что она рассматривается
как функция параметра $\theta$ и не обязана интегрироваться в единицу по
$\theta$.

Оценка максимального правдоподобия имеет вид
\[
    \widehat{\theta}_{\mathrm{MLE}}
    \in
    \argmax_{\theta}p(x\mid\theta),
\]
а максимум апостериорной плотности определяется формулой
\[
    \widehat{\theta}_{\mathrm{MAP}}
    \in
    \argmax_{\theta}
    \bigl\{
        p(x\mid\theta)\pi(\theta)
    \bigr\}.
\]
Эквивалентно,
\[
    \widehat{\theta}_{\mathrm{MAP}}
    \in
    \argmax_{\theta}
    \bigl\{
        \log p(x\mid\theta)+\log\pi(\theta)
    \bigr\}.
\]

\subsection{Пример применения теоремы Байеса}

Пусть:
\begin{itemize}
    \item доля объектов с некоторым свойством $D$ равна
    \[
        \Pp(D)=0{,}01;
    \]
    \item чувствительность теста равна
    \[
        \Pp(+\mid D)=0{,}95;
    \]
    \item специфичность теста равна
    \[
        \Pp(-\mid\overline{D})=0{,}90.
    \]
\end{itemize}

Тогда вероятность ложноположительного результата равна
\[
    \Pp(+\mid\overline{D})=1-0{,}90=0{,}10.
\]

Вероятностное дерево имеет следующий вид:

\begin{center}
\begin{forest}
for tree={
    draw,
    rounded corners,
    align=center,
    edge={->},
    l sep=13mm,
    s sep=12mm,
    font=\small
}
[{Случайно выбранный\\объект}
    [{\(D\)\\\(\Pp(D)=0{,}01\)}
        [{\(+\)\\\(\Pp(+\mid D)=0{,}95\)}]
        [{\(-\)\\\(\Pp(-\mid D)=0{,}05\)}]
    ]
    [{\(\overline{D}\)\\\(\Pp(\overline{D})=0{,}99\)}
        [{\(+\)\\\(\Pp(+\mid\overline{D})=0{,}10\)}]
        [{\(-\)\\\(\Pp(-\mid\overline{D})=0{,}90\)}]
    ]
]
\end{forest}
\end{center}

По формуле полной вероятности
\[
\begin{aligned}
    \Pp(+)
    &=
    \Pp(+\mid D)\Pp(D)
    +
    \Pp(+\mid\overline{D})\Pp(\overline{D})
    \\
    &=
    0{,}95\cdot 0{,}01
    +
    0{,}10\cdot 0{,}99
    \\
    &=
    0{,}0095+0{,}099
    =
    0{,}1085.
\end{aligned}
\]

По теореме Байеса
\[
\begin{aligned}
    \Pp(D\mid +)
    &=
    \frac{\Pp(+\mid D)\Pp(D)}{\Pp(+)}
    \\
    &=
    \frac{0{,}95\cdot 0{,}01}{0{,}1085}
    =
    \frac{0{,}0095}{0{,}1085}
    \approx 0{,}0876.
\end{aligned}
\]

Таким образом, при положительном результате теста искомое свойство
присутствует примерно с вероятностью $8{,}76\%$.

Результат удобно интерпретировать с помощью абсолютных частот. Среди
$10000$ объектов:
\begin{itemize}
    \item свойство имеют $100$ объектов, из них $95$ получат положительный
    результат;
    \item свойства не имеют $9900$ объектов, из них $990$ получат
    ложноположительный результат;
    \item всего положительный результат получат $95+990=1085$ объектов.
\end{itemize}
Следовательно,
\[
    \Pp(D\mid +)
    =
    \frac{95}{1085}
    \approx 0{,}0876.
\]

\begin{importantblock}{Основной вывод}
Высокая чувствительность теста сама по себе не гарантирует высокой
апостериорной вероятности. Необходимо учитывать базовую частоту
$\Pp(D)$ и вероятность ложноположительного результата.
\end{importantblock}

% ============================================================
\section{Байесовское правило принятия решений}
% ============================================================

\subsection{Функция потерь и апостериорный риск}

Пусть истинное состояние принадлежит множеству классов
\[
    Y\in\{1,\ldots,K\},
\]
а множество допустимых решений обозначено через $\mathcal{A}$.

Функция
\[
    L(a,k)
\]
задаёт потери, возникающие при выборе решения $a\in\mathcal{A}$, если
истинный класс равен $k$.

После наблюдения признаков $X=x$ определяется \emph{условный, или
апостериорный, риск}:
\[
    \rho(a\mid x)
    =
    \sum_{k=1}^{K}
    L(a,k)\Pp(Y=k\mid X=x).
\]

\begin{definition}
\emph{Байесовским решением} называется решение, минимизирующее
апостериорный риск:
\[
    a^*(x)
    \in
    \argmin_{a\in\mathcal{A}}
    \rho(a\mid x).
\]
\end{definition}

\subsection{Байесовский классификатор}

Рассмотрим функцию потерь $0$--$1$:
\[
    L(a,k)
    =
    \begin{cases}
        0, & a=k,\\
        1, & a\ne k.
    \end{cases}
\]

Тогда условный риск решения $a$ равен
\[
\begin{aligned}
    \rho(a\mid x)
    &=
    \sum_{k=1}^{K}
    \ind\{a\ne k\}\Pp(Y=k\mid X=x)
    \\
    &=
    1-\Pp(Y=a\mid X=x).
\end{aligned}
\]

Следовательно, необходимо выбирать класс с наибольшей апостериорной
вероятностью:
\[
    g^*(x)
    \in
    \argmax_{k\in\{1,\ldots,K\}}
    \Pp(Y=k\mid X=x).
\]

\begin{theorem}[Оптимальность байесовского классификатора]
Пусть качество классификатора $g$ измеряется вероятностью ошибки
\[
    \Risk(g)
    =
    \Pp(g(X)\ne Y).
\]
Тогда байесовский классификатор
\[
    g^*(x)
    \in
    \argmax_k \Pp(Y=k\mid X=x)
\]
минимизирует вероятность ошибки среди всех измеримых классификаторов:
\[
    \Risk(g^*)\le \Risk(g)
\]
для любого классификатора $g$.
\end{theorem}

\begin{proof}
Запишем риск классификатора через индикатор ошибки:
\[
    \Risk(g)
    =
    \Ee\bigl[\ind\{g(X)\ne Y\}\bigr].
\]
Применим формулу полного математического ожидания:
\[
\begin{aligned}
    \Risk(g)
    &=
    \Ee\left[
        \Ee\left[
            \ind\{g(X)\ne Y\}
            \mid X
        \right]
    \right].
\end{aligned}
\]

При фиксированном $X=x$ условная вероятность ошибки равна
\[
\begin{aligned}
    \Ee\left[
        \ind\{g(X)\ne Y\}
        \mid X=x
    \right]
    &=
    \Pp(g(x)\ne Y\mid X=x)
    \\
    &=
    1-\Pp(Y=g(x)\mid X=x).
\end{aligned}
\]

Чтобы минимизировать это выражение, необходимо максимизировать
$\Pp(Y=g(x)\mid X=x)$. Поэтому для почти всех $x$ оптимальным является
решение
\[
    g^*(x)
    \in
    \argmax_k\Pp(Y=k\mid X=x).
\]

Следовательно, для любого $g$ и почти всех $x$
\[
    \Pp(g^*(x)\ne Y\mid X=x)
    \le
    \Pp(g(x)\ne Y\mid X=x).
\]
Беря математическое ожидание по $X$, получаем
\[
    \Risk(g^*)\le \Risk(g).
\]
\end{proof}

Минимально достижимая вероятность ошибки называется
\emph{байесовской ошибкой}:
\[
    \Risk^*
    =
    \Ee\left[
        1-\max_k \Pp(Y=k\mid X)
    \right].
\]

\subsection{Различная стоимость ошибок}

Для бинарной классификации пусть:
\begin{itemize}
    \item $c_{10}$ --- стоимость решения $1$, когда истинный класс равен
    $0$;
    \item $c_{01}$ --- стоимость решения $0$, когда истинный класс равен
    $1$;
    \item правильные решения имеют нулевую стоимость.
\end{itemize}

Условный риск решения $1$:
\[
    \rho(1\mid x)
    =
    c_{10}\Pp(Y=0\mid X=x).
\]
Условный риск решения $0$:
\[
    \rho(0\mid x)
    =
    c_{01}\Pp(Y=1\mid X=x).
\]

Нужно выбрать класс $1$, если
\[
    c_{10}\Pp(Y=0\mid X=x)
    \le
    c_{01}\Pp(Y=1\mid X=x).
\]
Эквивалентно,
\[
    \frac{\Pp(Y=1\mid X=x)}
         {\Pp(Y=0\mid X=x)}
    \ge
    \frac{c_{10}}{c_{01}}.
\]

Таким образом, оптимальный порог зависит не только от вероятностей
классов, но и от стоимости различных ошибок.

% ============================================================
\section{Деревья решений в машинном обучении}
% ============================================================

\subsection{Постановка задачи}

Пусть дана обучающая выборка
\[
    \mathcal{D}
    =
    \{(x_i,y_i)\}_{i=1}^{n},
\]
где
\[
    x_i=(x_{i1},\ldots,x_{id})\in\mathcal{X}
\]
--- вектор признаков, а $y_i$ --- целевая переменная.

В задаче классификации
\[
    y_i\in\{1,\ldots,K\}.
\]
В задаче регрессии
\[
    y_i\in\Rr.
\]

\begin{definition}
\emph{Дерево решений} --- иерархическая модель, состоящая из:
\begin{itemize}
    \item корневой вершины;
    \item внутренних вершин, содержащих предикаты над признаками;
    \item рёбер, соответствующих результатам проверки предикатов;
    \item листьев, содержащих итоговые предсказания.
\end{itemize}
\end{definition}

Для числового признака типичный предикат имеет вид
\[
    x_j\le s.
\]
Он разбивает множество объектов на две части:
\[
    R_{\mathrm{left}}
    =
    \{x:x_j\le s\},
    \qquad
    R_{\mathrm{right}}
    =
    \{x:x_j>s\}.
\]

Для категориального признака могут использоваться предикаты вида
\[
    x_j\in C,
\]
где $C$ --- некоторое подмножество категорий.

\subsection{Геометрическая интерпретация}

Пусть листья дерева задают области
\[
    R_1,\ldots,R_M,
\]
которые образуют разбиение пространства признаков:
\[
    R_i\cap R_j=\varnothing,\quad i\ne j,
    \qquad
    \bigcup_{m=1}^{M}R_m=\mathcal{X}.
\]

В каждом листе предсказание постоянно. Поэтому дерево можно записать
как кусочно-постоянную функцию:
\[
    \widehat{f}(x)
    =
    \sum_{m=1}^{M}
    c_m\ind\{x\in R_m\}.
\]

В регрессии $c_m$ является числом. В классификации $c_m$ может быть:
\begin{itemize}
    \item меткой класса;
    \item вектором оценок вероятностей классов;
    \item вектором оценок стоимости решений.
\end{itemize}

Если используются предикаты $x_j\le s$, дерево формирует области,
границы которых параллельны координатным осям.

\subsection{Алгоритм получения предсказания}

\begin{algorithmblock}{Алгоритм прохождения объекта по дереву}
Для объекта $x$:
\begin{enumerate}[label=\arabic*.]
    \item Начать с корневой вершины.
    \item Проверить предикат, записанный в текущей вершине.
    \item Перейти по ребру, соответствующему результату проверки.
    \item Повторять шаги 2--3, пока не будет достигнут лист.
    \item Вернуть предсказание, хранящееся в листе:
    \begin{itemize}
        \item класс или вероятности классов для классификации;
        \item числовое значение для регрессии.
    \end{itemize}
\end{enumerate}
\end{algorithmblock}

Время получения предсказания пропорционально глубине листа, в который
попал объект.

\subsection{Выразительная способность}

\begin{theorem}
Любая булева функция
\[
    f:\{0,1\}^{d}\to\{0,1\}
\]
может быть представлена бинарным деревом решений глубины не более $d$.
\end{theorem}

\begin{proof}
Построим полное бинарное дерево глубины $d$. В корне проверим значение
$x_1$. На втором уровне проверим $x_2$, на третьем --- $x_3$ и так далее.
На уровне $d$ проверим значение $x_d$.

Каждый путь от корня к листу однозначно задаёт набор значений
\[
    (a_1,\ldots,a_d)\in\{0,1\}^{d}.
\]
Таких наборов существует $2^d$, поэтому листья полного дерева находятся
во взаимно однозначном соответствии со всеми возможными входами.

В листе, соответствующем набору $(a_1,\ldots,a_d)$, запишем значение
\[
    f(a_1,\ldots,a_d).
\]
Тогда для любого входа $x$ прохождение по дереву закончится в листе,
который содержит ровно $f(x)$. Следовательно, построенное дерево
представляет функцию $f$.
\end{proof}

Теорема показывает высокую выразительную способность деревьев. Однако
полное дерево может иметь $2^d$ листьев, поэтому на практике требуется
контроль сложности.

\subsection{Предсказание в листе классификационного дерева}

Пусть в вершину $t$ попало $n_t$ обучающих объектов. Обозначим через
\[
    n_{tk}
    =
    \sum_{i:x_i\in t}\ind\{y_i=k\}
\]
число объектов класса $k$. Эмпирическая доля класса $k$ равна
\[
    \widehat{p}_{tk}
    =
    \frac{n_{tk}}{n_t}.
\]

\begin{proposition}
При функции потерь $0$--$1$ оптимальным предсказанием в листе является
класс с наибольшим числом обучающих объектов:
\[
    \widehat{c}_t
    \in
    \argmax_k n_{tk}
    =
    \argmax_k\widehat{p}_{tk}.
\]
\end{proposition}

\begin{proof}
Если в листе всегда предсказывать класс $c$, то ошибочно
классифицированы все объекты, кроме объектов класса $c$. Число ошибок
равно
\[
    n_t-n_{tc}.
\]
Минимизация этой величины эквивалентна максимизации $n_{tc}$. Поэтому
оптимальным является любой класс, имеющий максимальную частоту.
\end{proof}

Оценку вероятности класса можно положить равной
\[
    \widehat{\Pp}(Y=k\mid x\in t)
    =
    \widehat{p}_{tk}.
\]

Для сглаживания может использоваться формула
\[
    \widehat{p}_{tk}^{\,(\lambda)}
    =
    \frac{n_{tk}+\lambda}
         {n_t+K\lambda},
    \qquad \lambda>0.
\]

\subsection{Предсказание в листе регрессионного дерева}

Пусть в лист $t$ попали значения
\[
    y_1,\ldots,y_{n_t}.
\]

\begin{theorem}
Константа, минимизирующая сумму квадратов ошибок
\[
    Q(c)
    =
    \sum_{i=1}^{n_t}(y_i-c)^2,
\]
равна среднему значению
\[
    \overline{y}_t
    =
    \frac{1}{n_t}\sum_{i=1}^{n_t}y_i.
\]
\end{theorem}

\begin{proof}
Для каждого $i$ запишем
\[
    y_i-c
    =
    (y_i-\overline{y}_t)+(\overline{y}_t-c).
\]
Тогда
\[
\begin{aligned}
    \sum_{i=1}^{n_t}(y_i-c)^2
    &=
    \sum_{i=1}^{n_t}
    \left[
        (y_i-\overline{y}_t)
        +
        (\overline{y}_t-c)
    \right]^2
    \\
    &=
    \sum_{i=1}^{n_t}(y_i-\overline{y}_t)^2
    +
    2(\overline{y}_t-c)
    \sum_{i=1}^{n_t}(y_i-\overline{y}_t)
    \\
    &\quad+
    n_t(\overline{y}_t-c)^2.
\end{aligned}
\]
По определению среднего
\[
    \sum_{i=1}^{n_t}(y_i-\overline{y}_t)=0.
\]
Следовательно,
\[
    \sum_{i=1}^{n_t}(y_i-c)^2
    =
    \sum_{i=1}^{n_t}(y_i-\overline{y}_t)^2
    +
    n_t(\overline{y}_t-c)^2.
\]
Первое слагаемое не зависит от $c$, а второе неотрицательно и обращается
в ноль тогда и только тогда, когда
\[
    c=\overline{y}_t.
\]
Следовательно, среднее значение является точкой минимума.
\end{proof}

% ============================================================
\section{Критерии качества разбиения}
% ============================================================

\subsection{Неоднородность классификационной вершины}

Критерий неоднородности должен принимать малое значение, когда почти все
объекты принадлежат одному классу, и большое значение, когда классы
представлены примерно одинаково.

Пусть
\[
    \widehat{p}_{tk}
    =
    \frac{n_{tk}}{n_t}.
\]

\subsubsection{Доля ошибок классификации}

\[
    I_{\mathrm{err}}(t)
    =
    1-\max_k\widehat{p}_{tk}.
\]

Эта величина совпадает с минимальной долей ошибок в вершине при
предсказании одного класса.

\subsubsection{Критерий Джини}

\[
    \Gini(t)
    =
    1-\sum_{k=1}^{K}\widehat{p}_{tk}^{\,2}.
\]

Для двух классов, если доля первого класса равна $p$, то
\[
    \Gini(p)
    =
    1-p^2-(1-p)^2
    =
    2p(1-p).
\]

Критерий Джини:
\begin{itemize}
    \item равен нулю в чистой вершине;
    \item максимален при равных долях классов;
    \item сильнее реагирует на изменение состава вершины, чем доля
    ошибок классификации.
\end{itemize}

\subsubsection{Энтропия}

\[
    \Entropy(t)
    =
    -\sum_{k=1}^{K}
    \widehat{p}_{tk}
    \log_2\widehat{p}_{tk},
\]
где по непрерывности принимается
\[
    0\log_2 0=0.
\]

Для двух классов
\[
    h(p)
    =
    -p\log_2p-(1-p)\log_2(1-p).
\]

Энтропия:
\begin{itemize}
    \item равна нулю в чистой вершине;
    \item максимальна при равномерном распределении классов;
    \item имеет информационно-теоретическую интерпретацию.
\end{itemize}

\subsection{Неоднородность регрессионной вершины}

Для регрессии обычно используется среднеквадратичное отклонение:
\[
    I_{\mathrm{reg}}(t)
    =
    \frac{1}{n_t}
    \sum_{i:x_i\in t}
    (y_i-\overline{y}_t)^2.
\]

Также можно работать с суммой квадратов ошибок:
\[
    \operatorname{SSE}(t)
    =
    \sum_{i:x_i\in t}
    (y_i-\overline{y}_t)^2.
\]

\subsection{Уменьшение неоднородности}

Пусть разбиение $s$ делит вершину $t$ на дочерние вершины
$t_1,\ldots,t_q$. Качество разбиения определяется уменьшением
неоднородности:
\[
    \Delta I(s,t)
    =
    I(t)
    -
    \sum_{r=1}^{q}
    \frac{n_{t_r}}{n_t}I(t_r).
\]

Выбирается разбиение
\[
    s^*
    \in
    \argmax_s \Delta I(s,t).
\]

Для бинарного разбиения
\[
    \Delta I(s,t)
    =
    I(t)
    -
    \frac{n_L}{n_t}I(t_L)
    -
    \frac{n_R}{n_t}I(t_R).
\]

Если в качестве неоднородности используется энтропия, величина
\[
    \Gain(s,t)
    =
    \Entropy(t)
    -
    \sum_{r=1}^{q}
    \frac{n_{t_r}}{n_t}\Entropy(t_r)
\]
называется \emph{информационным выигрышем}.

Если рассматривать номер дочерней ветви как дискретную случайную
величину $B$, то
\[
    \sum_{r=1}^{q}
    \frac{n_{t_r}}{n_t}\Entropy(t_r)
\]
является эмпирической условной энтропией класса при известной ветви.
Поэтому
\[
    \Gain(s,t)
    =
    \Entropy(Y)-\Entropy(Y\mid B),
\]
то есть информационный выигрыш совпадает с эмпирической взаимной
информацией между классом и результатом проверки предиката.

\subsection{Выбор порога для числового признака}

Пусть в вершине находятся различные значения числового признака
\[
    z_{(1)}<z_{(2)}<\ldots<z_{(r)}.
\]

Достаточно рассматривать пороги
\[
    s_\ell
    =
    \frac{z_{(\ell)}+z_{(\ell+1)}}{2},
    \qquad
    \ell=1,\ldots,r-1.
\]

Действительно, любой порог внутри интервала
\[
    (z_{(\ell)},z_{(\ell+1)})
\]
создаёт одно и то же разбиение обучающих объектов. Поэтому проверка
других порогов внутри этого интервала не изменит значение критерия.

\subsection{Категориальные признаки}

Пусть категориальный признак принимает значения из множества
\[
    \{a_1,\ldots,a_q\}.
\]

Возможны два подхода:
\begin{enumerate}
    \item \textbf{Многоветвевое разбиение:}
    создаётся отдельная ветвь для каждого значения признака.
    Такой подход характерен для ID3 и C4.5.
    \item \textbf{Бинарное разбиение:}
    выбирается подмножество категорий $C$, после чего проверяется
    предикат
    \[
        x_j\in C.
    \]
    Такой подход характерен для CART.
\end{enumerate}

Полный перебор бинарных разбиений категорий требует рассмотрения
экспоненциального числа подмножеств. Поэтому при большом числе категорий
применяются эвристики, упорядочивание категорий или предварительное
кодирование.

% ============================================================
\section{Построение дерева}
% ============================================================

\subsection{Жадный рекурсивный алгоритм}

Глобальный поиск дерева оптимальной структуры является сложной
комбинаторной задачей. На практике дерево строится жадно: в каждой
вершине выбирается лучшее локальное разбиение.

\begin{algorithmblock}{Общий алгоритм построения дерева решений}
Вход:
\begin{itemize}
    \item обучающие объекты, попавшие в текущую вершину $t$;
    \item набор допустимых признаков и предикатов;
    \item критерий неоднородности;
    \item параметры остановки.
\end{itemize}

Шаги:
\begin{enumerate}[label=\arabic*.]
    \item Вычислить предсказание, которое будет храниться в вершине, если
    она станет листом:
    \begin{itemize}
        \item класс большинства для классификации;
        \item среднее значение целевой переменной для регрессии.
    \end{itemize}

    \item Проверить критерии остановки. Если дальнейшее разбиение
    запрещено, вернуть лист.

    \item Для каждого допустимого признака перебрать допустимые
    разбиения.

    \item Для каждого разбиения вычислить уменьшение неоднородности
    \[
        \Delta I(s,t).
    \]

    \item Выбрать лучшее разбиение
    \[
        s^*
        \in
        \argmax_s\Delta I(s,t).
    \]

    \item Если улучшение недостаточно или разбиение создаёт недопустимые
    дочерние вершины, вернуть лист.

    \item Создать внутреннюю вершину с предикатом $s^*$.

    \item Разделить обучающие объекты между дочерними вершинами.

    \item Рекурсивно построить поддерево для каждой дочерней вершины.

    \item Вернуть построенное дерево.
\end{enumerate}
\end{algorithmblock}

Поскольку решение на каждом шаге принимается локально, жадный алгоритм
не гарантирует нахождения глобально оптимального дерева.

\subsection{Критерии остановки}

Построение ветви может быть остановлено, если выполнено хотя бы одно из
условий:
\begin{itemize}
    \item все объекты в вершине принадлежат одному классу;
    \item достигнута максимальная глубина;
    \item число объектов меньше заданного порога;
    \item дальнейшее разбиение создаст слишком маленький лист;
    \item уменьшение неоднородности меньше заданного порога;
    \item исчерпаны допустимые признаки или разбиения;
    \item достигнуто максимальное число листьев.
\end{itemize}

Такие ограничения называются \emph{предварительной обрезкой}, или
\emph{предобрезкой} дерева.

% ============================================================
\section{Алгоритмы ID3, C4.5 и CART}
% ============================================================

\subsection{Алгоритм ID3}

Алгоритм ID3 предназначен прежде всего для классификации объектов,
описанных категориальными признаками. Для выбора признака используется
информационный выигрыш.

\begin{algorithmblock}{Алгоритм ID3}
Пусть $S$ --- множество объектов в текущей вершине, а $A$ --- множество
доступных признаков.

\begin{enumerate}[label=\arabic*.]
    \item Если все объекты $S$ принадлежат одному классу, создать лист с
    этим классом.

    \item Если множество доступных признаков пусто, создать лист с
    классом большинства в $S$.

    \item Для каждого признака $a\in A$ вычислить информационный выигрыш
    \[
        \Gain(S,a)
        =
        \Entropy(S)
        -
        \sum_{v\in\operatorname{Values}(a)}
        \frac{|S_v|}{|S|}
        \Entropy(S_v),
    \]
    где
    \[
        S_v=\{x\in S:a(x)=v\}.
    \]

    \item Выбрать признак
    \[
        a^*
        \in
        \argmax_{a\in A}\Gain(S,a).
    \]

    \item Создать внутреннюю вершину, проверяющую значение $a^*$.

    \item Для каждого значения $v$ признака $a^*$:
    \begin{enumerate}[label=\alph*)]
        \item создать ветвь $a^*=v$;
        \item рекурсивно построить дерево на множестве $S_v$;
        \item удалить $a^*$ из набора доступных признаков на этой ветви.
    \end{enumerate}

    \item Если для некоторого значения признака обучающих объектов нет,
    создать лист с классом большинства родительской вершины.
\end{enumerate}
\end{algorithmblock}

Недостаток информационного выигрыша состоит в предпочтении признаков с
большим числом различных значений. Например, уникальный идентификатор
может идеально разделить обучающую выборку, но не иметь предсказательной
ценности на новых объектах.

\subsection{Алгоритм C4.5}

C4.5 развивает идеи ID3. Одним из основных отличий является
использование \emph{отношения информационного выигрыша}.

Для разбиения на подмножества $S_1,\ldots,S_q$ определяется
\[
    \operatorname{SplitInfo}(S,a)
    =
    -
    \sum_{r=1}^{q}
    \frac{|S_r|}{|S|}
    \log_2\frac{|S_r|}{|S|}.
\]

Отношение выигрыша:
\[
    \operatorname{GainRatio}(S,a)
    =
    \frac{\Gain(S,a)}
         {\operatorname{SplitInfo}(S,a)}.
\]

В упрощённом изложении выбирается разбиение с наибольшим отношением
выигрыша среди разбиений, обеспечивающих содержательный информационный
выигрыш.

C4.5 также поддерживает:
\begin{itemize}
    \item числовые признаки и пороговые разбиения;
    \item пропущенные значения;
    \item многоклассовую классификацию;
    \item постобрезку дерева;
    \item преобразование дерева в набор правил.
\end{itemize}

\subsection{Алгоритм CART}

Название CART означает \emph{Classification and Regression Trees}.
Основные особенности:
\begin{itemize}
    \item только бинарные разбиения;
    \item поддержка классификации и регрессии;
    \item для классификации часто используется критерий Джини;
    \item для регрессии используется уменьшение суммы квадратов ошибок;
    \item применяется постобрезка по критерию ``ошибка плюс сложность''.
\end{itemize}

Для классификации выбирается разбиение, максимизирующее
\[
    \Delta\Gini
    =
    \Gini(t)
    -
    \frac{n_L}{n_t}\Gini(t_L)
    -
    \frac{n_R}{n_t}\Gini(t_R).
\]

Для регрессии выбирается разбиение, максимизирующее
\[
    \Delta\operatorname{SSE}
    =
    \operatorname{SSE}(t)
    -
    \operatorname{SSE}(t_L)
    -
    \operatorname{SSE}(t_R).
\]

\subsection{Сравнение алгоритмов}

\begin{center}
\small
\begin{tabularx}{\textwidth}{
    >{\raggedright\arraybackslash}p{2.2cm}
    >{\raggedright\arraybackslash}X
    >{\raggedright\arraybackslash}X
    >{\raggedright\arraybackslash}X
}
\toprule
\textbf{Свойство}
&
\textbf{ID3}
&
\textbf{C4.5}
&
\textbf{CART}
\\
\midrule

Задача
&
Классификация
&
Классификация
&
Классификация и регрессия
\\

Тип ветвления
&
Обычно многоветвевое
&
Многоветвевое для категориальных и пороговое для числовых признаков
&
Бинарное
\\

Основной критерий
&
Информационный выигрыш
&
Отношение информационного выигрыша
&
Джини для классификации, сумма квадратов для регрессии
\\

Числовые признаки
&
В исходной форме ограниченно
&
Поддерживаются
&
Поддерживаются
\\

Обрезка
&
В исходном алгоритме отсутствует развитая постобрезка
&
Постобрезка на основе оценки ошибки
&
Обрезка ``ошибка плюс сложность''
\\

Пропущенные значения
&
Ограниченная поддержка
&
Поддерживаются
&
Могут использоваться замещающие разбиения
\\

\bottomrule
\end{tabularx}
\end{center}

% ============================================================
\section{Минимальный пример классификационного дерева}
% ============================================================

Рассмотрим выборку с двумя бинарными признаками $x_1,x_2$ и классами
$+$ и $-$:

\begin{center}
\begin{tabular}{cccc}
\toprule
$i$ & $x_1$ & $x_2$ & $y$\\
\midrule
1 & 1 & 1 & $+$\\
2 & 1 & 0 & $+$\\
3 & 1 & 1 & $+$\\
4 & 1 & 0 & $+$\\
5 & 0 & 1 & $+$\\
6 & 0 & 0 & $-$\\
7 & 0 & 0 & $-$\\
8 & 0 & 0 & $-$\\
\bottomrule
\end{tabular}
\end{center}

В корневой вершине находится пять объектов класса $+$ и три объекта
класса $-$. Энтропия равна
\[
\begin{aligned}
    \Entropy(S)
    &=
    -\frac{5}{8}\log_2\frac{5}{8}
    -\frac{3}{8}\log_2\frac{3}{8}
    \\
    &\approx 0{,}9544.
\end{aligned}
\]

\subsection*{Разбиение по признаку $x_1$}

При $x_1=1$ получаем четыре объекта класса $+$:
\[
    \Entropy(S_{x_1=1})=0.
\]

При $x_1=0$ получаем один объект класса $+$ и три объекта класса $-$:
\[
\begin{aligned}
    \Entropy(S_{x_1=0})
    &=
    -\frac14\log_2\frac14
    -\frac34\log_2\frac34
    \\
    &\approx 0{,}8113.
\end{aligned}
\]

Средняя энтропия после разбиения:
\[
    \frac48\cdot 0
    +
    \frac48\cdot 0{,}8113
    \approx 0{,}4056.
\]

Информационный выигрыш:
\[
    \Gain(S,x_1)
    =
    0{,}9544-0{,}4056
    \approx 0{,}5488.
\]

\subsection*{Разбиение по признаку $x_2$}

При $x_2=1$ находятся три объекта класса $+$, поэтому энтропия равна
нулю.

При $x_2=0$ находятся два объекта класса $+$ и три объекта класса $-$:
\[
\begin{aligned}
    \Entropy(S_{x_2=0})
    &=
    -\frac25\log_2\frac25
    -\frac35\log_2\frac35
    \\
    &\approx 0{,}9710.
\end{aligned}
\]

Средняя энтропия после разбиения:
\[
    \frac38\cdot 0
    +
    \frac58\cdot 0{,}9710
    \approx 0{,}6068.
\]

Информационный выигрыш:
\[
    \Gain(S,x_2)
    =
    0{,}9544-0{,}6068
    \approx 0{,}3476.
\]

Поскольку
\[
    \Gain(S,x_1)>\Gain(S,x_2),
\]
в корне выбирается признак $x_1$.

В ветви $x_1=0$ признак $x_2$ полностью разделяет классы. Итоговое
дерево:

\begin{center}
\begin{forest}
for tree={
    draw,
    rounded corners,
    align=center,
    edge={->},
    l sep=14mm,
    s sep=15mm,
    font=\small
}
[{\(x_1=1?\)}
    [{Да: класс \(+\)}]
    [{Нет: \(x_2=1?\)}
        [{Да: класс \(+\)}]
        [{Нет: класс \(-\)}]
    ]
]
\end{forest}
\end{center}

Дерево задаёт правило
\[
    \widehat{y}
    =
    \begin{cases}
        +, & x_1=1,\\
        +, & x_1=0,\ x_2=1,\\
        -, & x_1=0,\ x_2=0.
    \end{cases}
\]

Эквивалентная логическая запись:
\[
    \widehat{y}=+
    \quad\Longleftrightarrow\quad
    (x_1=1)\lor(x_1=0\land x_2=1).
\]

Таким образом, каждый путь к листу соответствует конъюнкции условий, а
объединение листьев одного класса --- дизъюнкции таких конъюнкций.

% ============================================================
\section{Минимальный пример регрессионного дерева}
% ============================================================

Рассмотрим выборку
\[
\begin{array}{c|cccc}
    x & 1 & 2 & 3 & 4\\
    \hline
    y & 1 & 1 & 4 & 4
\end{array}.
\]

Если не выполнять разбиение, предсказанием будет среднее
\[
    \overline{y}
    =
    \frac{1+1+4+4}{4}
    =
    2{,}5.
\]

Сумма квадратов ошибок в корне:
\[
\begin{aligned}
    \operatorname{SSE}(t)
    &=
    (1-2{,}5)^2
    +(1-2{,}5)^2
    +(4-2{,}5)^2
    +(4-2{,}5)^2
    \\
    &=9.
\end{aligned}
\]

Рассмотрим порог
\[
    x\le 2{,}5.
\]

В левой вершине находятся значения $1,1$, поэтому
\[
    \overline{y}_L=1,
    \qquad
    \operatorname{SSE}(t_L)=0.
\]

В правой вершине находятся значения $4,4$, поэтому
\[
    \overline{y}_R=4,
    \qquad
    \operatorname{SSE}(t_R)=0.
\]

Уменьшение ошибки:
\[
    \Delta\operatorname{SSE}
    =
    9-0-0
    =
    9.
\]

Итоговое дерево:

\begin{center}
\begin{forest}
for tree={
    draw,
    rounded corners,
    align=center,
    edge={->},
    l sep=14mm,
    s sep=17mm,
    font=\small
}
[{\(x\le 2{,}5?\)}
    [{Да: \(\widehat{y}=1\)}]
    [{Нет: \(\widehat{y}=4\)}]
]
\end{forest}
\end{center}

Соответствующая функция имеет вид
\[
    \widehat{f}(x)
    =
    \begin{cases}
        1, & x\le 2{,}5,\\
        4, & x>2{,}5.
    \end{cases}
\]

% ============================================================
\section{Переобучение и обрезка дерева}
% ============================================================

\subsection{Причина переобучения}

Если разрешить неограниченный рост дерева, можно получить листья,
содержащие по одному объекту. Такое дерево будет иметь очень малую
ошибку на обучающей выборке, но может плохо работать на новых данных.

Глубокое дерево обычно характеризуется:
\begin{itemize}
    \item низким смещением;
    \item высокой дисперсией;
    \item чувствительностью к небольшим изменениям обучающей выборки;
    \item сложной структурой и большим числом листьев.
\end{itemize}

Контроль сложности выполняется с помощью:
\begin{itemize}
    \item предварительной обрезки;
    \item постобрезки уже построенного дерева.
\end{itemize}

\subsection{Предварительная обрезка}

К основным параметрам предварительной обрезки относятся:
\begin{itemize}
    \item максимальная глубина дерева;
    \item минимальное число объектов для разбиения вершины;
    \item минимальное число объектов в листе;
    \item минимальное уменьшение неоднородности;
    \item максимальное число листьев.
\end{itemize}

Недостаток предобрезки состоит в том, что локально слабое разбиение может
оказаться необходимым для последующего сильного разбиения. Слишком
ранняя остановка способна привести к недообучению.

\subsection{Обрезка ``ошибка плюс сложность''}

Пусть $T$ --- дерево, а $\mathcal{L}(T)$ --- множество его листьев.
Рассмотрим функционал
\[
    \Risk_\alpha(T)
    =
    \Risk(T)
    +
    \alpha|\mathcal{L}(T)|,
    \qquad \alpha\ge 0.
\]

Здесь:
\begin{itemize}
    \item $\Risk(T)$ --- эмпирическая ошибка дерева;
    \item $|\mathcal{L}(T)|$ --- число листьев;
    \item $\alpha$ --- штраф за сложность.
\end{itemize}

При $\alpha=0$ предпочтение отдаётся дереву с минимальной обучающей
ошибкой. При увеличении $\alpha$ более простые деревья становятся
предпочтительнее.

Для классификации можно использовать
\[
    \Risk(T)
    =
    \sum_{t\in\mathcal{L}(T)}
    \frac{n_t}{n}
    \left(
        1-\max_k\widehat{p}_{tk}
    \right).
\]

Для регрессии:
\[
    \Risk(T)
    =
    \frac{1}{n}
    \sum_{t\in\mathcal{L}(T)}
    \sum_{i:x_i\in t}
    (y_i-\overline{y}_t)^2.
\]

\subsection{Алгоритм weakest-link pruning}

Пусть $T_t$ --- поддерево, корнем которого является внутренняя вершина
$t$. Обозначим:
\begin{itemize}
    \item через $\Risk(T_t)$ ошибку всего поддерева;
    \item через $\Risk(t)$ ошибку, если заменить всё поддерево одним
    листом;
    \item через $|\mathcal{L}(T_t)|$ число листьев поддерева.
\end{itemize}

Критическое значение штрафа для вершины $t$:
\[
    \alpha_t
    =
    \frac{
        \Risk(t)-\Risk(T_t)
    }{
        |\mathcal{L}(T_t)|-1
    }.
\]

При $\alpha\ge\alpha_t$ замена поддерева $T_t$ одним листом становится
не хуже с точки зрения функционала $\Risk_\alpha$.

\begin{algorithmblock}{Постобрезка CART}
\begin{enumerate}[label=\arabic*.]
    \item Построить достаточно большое дерево $T_0$.

    \item Для каждой внутренней вершины $t$ вычислить
    \[
        \alpha_t
        =
        \frac{
            \Risk(t)-\Risk(T_t)
        }{
            |\mathcal{L}(T_t)|-1
        }.
    \]

    \item Найти минимальное значение $\alpha_t$.

    \item Удалить поддерево или поддеревья, соответствующие минимальному
    значению, заменив их листьями.

    \item Повторять шаги 2--4, пока дерево не сократится до одного листа.

    \item Получить последовательность вложенных деревьев
    \[
        T_0\supset T_1\supset\ldots\supset T_M.
    \]

    \item С помощью валидационной выборки или кросс-валидации выбрать
    дерево с наименьшей оценкой ошибки.
\end{enumerate}
\end{algorithmblock}

\begin{exampleblock}{Минимальный пример обрезки}
Пусть поддерево имеет три листа и ошибку
\[
    \Risk(T_t)=0{,}08.
\]
Если заменить его одним листом, ошибка станет
\[
    \Risk(t)=0{,}12.
\]
Тогда
\[
    \alpha_t
    =
    \frac{0{,}12-0{,}08}{3-1}
    =
    0{,}02.
\]

Для поддерева значение функционала равно
\[
    0{,}08+3\alpha,
\]
а для одного листа
\[
    0{,}12+\alpha.
\]
При $\alpha\ge 0{,}02$ один лист не хуже трёхлистного поддерева.
\end{exampleblock}

% ============================================================
\section{Практические аспекты деревьев решений}
% ============================================================

\subsection{Пропущенные значения}

Возможные способы обработки пропусков:
\begin{itemize}
    \item предварительное заполнение пропущенных значений;
    \item выделение пропуска в отдельную категорию;
    \item направление объекта в наиболее крупную дочернюю вершину;
    \item вероятностное распределение объекта между ветвями;
    \item использование замещающих разбиений.
\end{itemize}

\emph{Замещающее разбиение} --- предикат по другому признаку, который
наиболее точно воспроизводит результат основного разбиения на объектах,
где оба признака известны.

\subsection{Несбалансированные классы}

Если один класс встречается значительно реже другого, простая доля
правильно классифицированных объектов может быть неинформативна.

Возможные способы учёта дисбаланса:
\begin{itemize}
    \item назначение весов классам;
    \item изменение функции потерь;
    \item изменение порога решения;
    \item увеличение или уменьшение числа объектов отдельных классов;
    \item использование специальных метрик качества.
\end{itemize}

При весах объектов $w_i$ доля класса в вершине определяется как
\[
    \widehat{p}_{tk}^{(w)}
    =
    \frac{
        \displaystyle
        \sum_{i:x_i\in t}
        w_i\ind\{y_i=k\}
    }{
        \displaystyle
        \sum_{i:x_i\in t}w_i
    }.
\]

\subsection{Оценка важности признаков}

Важность признака может оцениваться суммарным уменьшением
неоднородности во всех вершинах, где используется этот признак:
\[
    \operatorname{Imp}(j)
    =
    \sum_{t:j(t)=j}
    \frac{n_t}{n}\Delta I(t).
\]

После нормировки:
\[
    \widetilde{\operatorname{Imp}}(j)
    =
    \frac{\operatorname{Imp}(j)}
         {\sum_{\ell=1}^{d}\operatorname{Imp}(\ell)}.
\]

Такая оценка может быть смещена в пользу:
\begin{itemize}
    \item числовых признаков с большим числом возможных порогов;
    \item категориальных признаков с большим числом категорий.
\end{itemize}

Альтернативой является перестановочная важность: значения одного
признака случайно перемешиваются, после чего измеряется ухудшение
качества модели.

\subsection{Интерпретация дерева}

Каждый путь от корня к листу представляет правило вида
\[
    \text{условие}_1
    \land
    \text{условие}_2
    \land\ldots\land
    \text{условие}_r
    \quad\Longrightarrow\quad
    \text{предсказание}.
\]

Для классификации в листе полезно хранить:
\begin{itemize}
    \item число обучающих объектов;
    \item распределение классов;
    \item прогнозируемый класс;
    \item оценку вероятности класса;
    \item значение эмпирического риска.
\end{itemize}

Интерпретируемость ухудшается при:
\begin{itemize}
    \item большой глубине;
    \item большом числе листьев;
    \item повторном использовании одних и тех же признаков;
    \item наличии сильно коррелированных признаков.
\end{itemize}

\subsection{Оценивание качества}

Данные обычно разделяются на:
\begin{itemize}
    \item обучающую выборку --- для построения дерева;
    \item валидационную выборку --- для выбора глубины, параметров
    остановки и степени обрезки;
    \item тестовую выборку --- для итоговой оценки качества.
\end{itemize}

Для классификации применяются:
\begin{itemize}
    \item accuracy;
    \item precision и recall;
    \item $F_1$-мера;
    \item ROC-AUC;
    \item PR-AUC;
    \item логарифмическая функция потерь;
    \item матрица ошибок.
\end{itemize}

Для регрессии:
\begin{itemize}
    \item среднеквадратичная ошибка;
    \item корень из среднеквадратичной ошибки;
    \item средняя абсолютная ошибка;
    \item коэффициент детерминации.
\end{itemize}

При выборе параметров дерева необходимо использовать кросс-валидацию
или отдельную валидационную выборку. Выбор параметров по тестовой
выборке приводит к смещённой оценке качества.

\subsection{Преимущества деревьев решений}

\begin{itemize}
    \item Наглядность и интерпретируемость небольших деревьев.
    \item Возможность работы с нелинейными зависимостями.
    \item Автоматический учёт взаимодействий признаков.
    \item Отсутствие необходимости стандартизовать числовые признаки.
    \item Возможность работать как с классификацией, так и с регрессией.
    \item Возможность использовать числовые и категориальные признаки.
    \item Быстрое получение предсказаний.
    \item Возможность извлекать логические правила.
\end{itemize}

\subsection{Недостатки деревьев решений}

\begin{itemize}
    \item Склонность глубоких деревьев к переобучению.
    \item Высокая чувствительность структуры дерева к изменениям данных.
    \item Жадное обучение не гарантирует глобально оптимальную структуру.
    \item Кусочно-постоянные предсказания в регрессии.
    \item Неустойчивая экстраполяция за пределы обучающего диапазона.
    \item Предпочтение некоторых критериев к признакам с большим числом
    возможных разбиений.
    \item Сложность представления наклонных или гладких границ решений
    одним небольшим деревом.
\end{itemize}

Для уменьшения дисперсии используются ансамбли деревьев:
\begin{itemize}
    \item бэггинг;
    \item случайный лес;
    \item градиентный бустинг.
\end{itemize}
Однако ансамбли обычно менее интерпретируемы, чем отдельное небольшое
дерево.

% ============================================================
\section{Деревья в теории принятия решений}
% ============================================================

Следует различать:
\begin{enumerate}
    \item \emph{обучаемое дерево решений} --- статистическую модель,
    строящуюся по обучающей выборке;
    \item \emph{дерево принятия решений} --- графическое представление
    последовательности решений, случайных событий и выигрышей.
\end{enumerate}

В дереве принятия решений используются:
\begin{itemize}
    \item вершины решения, в которых выбирается одно из действий;
    \item вероятностные вершины, в которых исход определяется случайным
    событием;
    \item терминальные вершины, содержащие итоговый выигрыш или потери.
\end{itemize}

\subsection{Метод обратной свёртки}

\begin{algorithmblock}{Вычисление оптимальной стратегии по дереву}
\begin{enumerate}[label=\arabic*.]
    \item Задать последовательность наблюдений, решений и случайных
    событий.

    \item Назначить вероятности ветвям, выходящим из вероятностных
    вершин.

    \item Если получено новое наблюдение, обновить вероятности гипотез по
    теореме Байеса.

    \item Записать выигрыши или потери в терминальных вершинах.

    \item Двигаться от терминальных вершин к корню:
    \begin{itemize}
        \item в вероятностной вершине вычислять математическое ожидание;
        \item в вершине решения выбирать действие с максимальным
        ожидаемым выигрышем или минимальными ожидаемыми потерями.
    \end{itemize}

    \item Сохранить выбранные действия. Они образуют оптимальную
    стратегию.
\end{enumerate}
\end{algorithmblock}

Если из вероятностной вершины выходят ветви с вероятностями $p_i$ и
значениями $V_i$, то значение вершины равно
\[
    V=\sum_i p_iV_i.
\]

Если из вершины решения доступны действия со значениями $V_1,\ldots,V_m$,
то при максимизации полезности
\[
    V=\max_i V_i.
\]

\subsection{Пример совместного применения дерева и теоремы Байеса}

Пусть механизм неисправен с априорной вероятностью
\[
    \Pp(F)=0{,}1.
\]
Датчик имеет характеристики
\[
    \Pp(+\mid F)=0{,}8,
    \qquad
    \Pp(+\mid\overline{F})=0{,}2.
\]

После положительного сигнала
\[
\begin{aligned}
    \Pp(F\mid +)
    &=
    \frac{0{,}8\cdot 0{,}1}
    {0{,}8\cdot 0{,}1+0{,}2\cdot 0{,}9}
    \\
    &=
    \frac{0{,}08}{0{,}26}
    =
    \frac{4}{13}
    \approx 0{,}3077.
\end{aligned}
\]

После отрицательного сигнала
\[
\begin{aligned}
    \Pp(F\mid -)
    &=
    \frac{0{,}2\cdot 0{,}1}
    {0{,}2\cdot 0{,}1+0{,}8\cdot 0{,}9}
    \\
    &=
    \frac{0{,}02}{0{,}74}
    =
    \frac{1}{37}
    \approx 0{,}0270.
\end{aligned}
\]

Пусть:
\begin{itemize}
    \item профилактический ремонт имеет стоимость $3$;
    \item отказ неисправного механизма без ремонта имеет стоимость $10$;
    \item если механизм исправен и ремонт не выполняется, потери равны
    нулю.
\end{itemize}

После положительного сигнала ожидаемые потери при отказе от ремонта:
\[
    10\Pp(F\mid +)
    =
    \frac{40}{13}
    \approx 3{,}0769.
\]
Поскольку
\[
    3<3{,}0769,
\]
после положительного сигнала следует выполнить ремонт.

После отрицательного сигнала ожидаемые потери при отказе от ремонта:
\[
    10\Pp(F\mid -)
    =
    \frac{10}{37}
    \approx 0{,}2703.
\]
Поскольку
\[
    0{,}2703<3,
\]
после отрицательного сигнала ремонт выполнять не следует.

Вероятности сигналов:
\[
    \Pp(+)=0{,}26,
    \qquad
    \Pp(-)=0{,}74.
\]

Ожидаемые потери оптимальной стратегии с бесплатным датчиком:
\[
\begin{aligned}
    \Ee[L\mid\text{датчик}]
    &=
    0{,}26\cdot 3
    +
    0{,}74\cdot\frac{10}{37}
    \\
    &=0{,}78+0{,}20
    =0{,}98.
\end{aligned}
\]

Без датчика:
\[
    \Ee[L\mid\text{не ремонтировать}]
    =
    10\cdot 0{,}1
    =
    1,
\]
а немедленный ремонт имеет стоимость $3$. Поэтому без датчика
оптимально не выполнять ремонт, получая ожидаемые потери $1$.

Ценность информации датчика:
\[
    1-0{,}98=0{,}02.
\]
Следовательно, использование датчика экономически оправдано, если его
стоимость меньше $0{,}02$ в выбранных единицах потерь.

% ============================================================
\section{Связь теоремы Байеса и обучаемых деревьев}
% ============================================================

Теорема Байеса и деревья решений относятся к разным уровням описания
задачи классификации.

Теорема Байеса определяет оптимальное решение, если известны истинные
апостериорные вероятности:
\[
    g^*(x)
    \in
    \argmax_k\Pp(Y=k\mid X=x).
\]

На практике распределение
\[
    \Pp(Y=k\mid X=x)
\]
неизвестно. Дерево решений оценивает его локально по объектам,
попавшим в соответствующий лист:
\[
    \widehat{\Pp}(Y=k\mid X=x)
    =
    \frac{n_{tk}}{n_t},
    \qquad x\in t.
\]

После этого можно применить байесовское правило решения:
\[
    \widehat{g}(x)
    \in
    \argmin_a
    \sum_{k=1}^{K}
    L(a,k)
    \widehat{\Pp}(Y=k\mid X=x).
\]

При функции потерь $0$--$1$ получаем класс большинства:
\[
    \widehat{g}(x)
    \in
    \argmax_k
    \widehat{\Pp}(Y=k\mid X=x).
\]

Таким образом:
\begin{itemize}
    \item теорема Байеса задаёт правило обновления вероятностей;
    \item байесовская теория решений задаёт оптимальное действие при
    известных вероятностях;
    \item дерево решений строит разбиение пространства признаков и
    оценивает распределение классов внутри каждой области.
\end{itemize}

% ============================================================
\section{Основные формулы}
% ============================================================

\begin{importantblock}{Условная вероятность}
\[
    \Pp(A\mid B)
    =
    \frac{\Pp(A\cap B)}{\Pp(B)}.
\]
\end{importantblock}

\begin{importantblock}{Формула полной вероятности}
\[
    \Pp(A)
    =
    \sum_{i=1}^{m}
    \Pp(A\mid H_i)\Pp(H_i).
\]
\end{importantblock}

\begin{importantblock}{Теорема Байеса}
\[
    \Pp(H_k\mid A)
    =
    \frac{
        \Pp(A\mid H_k)\Pp(H_k)
    }{
        \displaystyle
        \sum_{i=1}^{m}
        \Pp(A\mid H_i)\Pp(H_i)
    }.
\]
\end{importantblock}

\begin{importantblock}{Байесовский классификатор}
\[
    g^*(x)
    \in
    \argmax_k\Pp(Y=k\mid X=x).
\]
\end{importantblock}

\begin{importantblock}{Энтропия}
\[
    \Entropy(t)
    =
    -\sum_{k=1}^{K}
    \widehat{p}_{tk}\log_2\widehat{p}_{tk}.
\]
\end{importantblock}

\begin{importantblock}{Критерий Джини}
\[
    \Gini(t)
    =
    1-\sum_{k=1}^{K}\widehat{p}_{tk}^{\,2}.
\]
\end{importantblock}

\begin{importantblock}{Уменьшение неоднородности}
\[
    \Delta I(s,t)
    =
    I(t)
    -
    \sum_{r}
    \frac{n_{t_r}}{n_t}I(t_r).
\]
\end{importantblock}

\begin{importantblock}{Регрессионное дерево}
\[
    \widehat{y}_t
    =
    \frac{1}{n_t}
    \sum_{i:x_i\in t}y_i.
\]
\end{importantblock}

\begin{importantblock}{Обрезка ``ошибка плюс сложность''}
\[
    \Risk_\alpha(T)
    =
    \Risk(T)
    +
    \alpha|\mathcal{L}(T)|.
\]
\end{importantblock}

\end{document}
